\documentclass[10pt]{article}

% Specify the margins and text size.
\setlength{\textwidth}{6.4in}
\setlength{\textheight}{9.5in}
\setlength{\oddsidemargin}{0pt}
\setlength{\evensidemargin}{0pt}
\setlength{\topmargin}{0pt}
\setlength{\hoffset}{.05in}
\setlength{\voffset}{-1in}

\setlength{\parskip}{5pt}
\setlength{\parindent}{0pt}

% Load some fonts and symbol packages
\usepackage{latexsym}
\usepackage{pifont}       % contains 'star' symbol for counterinsurgency handbook title
%\usepackage{yfonts} 
\usepackage{amsmath}
\usepackage{amsfonts}

\usepackage{graphicx}     % actually, this is loaded by pstricks
\usepackage[T1]{fontenc}
\usepackage{ifthen}
\usepackage{pstricks,pst-grad,pst-text,pst-node,multido,pst-plot,calc,pst-3dplot}
%\usepackage[all]{xy}
%\usepackage{animate}

% The hyperref package inserts links into the pdf.
\definecolor{MyLinkColor}{rgb}{.1,.2,1}
\definecolor{MyCiteColor}{rgb}{.1,1,.2}
\definecolor{MyURLColor}{rgb}{.4,.4,.4}
\usepackage[backref=true,pagebackref=false,hyperindex,colorlinks=true,
  linkcolor=MyLinkColor,urlcolor=MyURLColor]{hyperref}


% The tweaklist package is something I found on the web.  It provides a simple interface
% for making changes to spacing used in the itemize and enumerate environments.  Comment
% this out if you don't care to use tweaklist.
\usepackage{tweaklist}
\renewcommand{\itemhook}{\setlength{\parskip}{2pt}\setlength{\parsep}%
{1pt}\setlength{\topsep}{0pt}\setlength{\itemsep}{0pt}}

\newcommand{\U}{\underline{\hspace{5pt}}}

\usepackage{listings}
\newcommand{\Z}{\hphantom{0}}

\begin{document}
\pagestyle{empty}
\lstset{language=R, showspaces=false, showstringspaces=false}
\href{http://www.su.edu}{\includegraphics[height=1.75cm]{sulogo.eps}}
\vspace{-1.79cm}

{{\ }\hfill\small
\begin{tabular}{cl}
& Math 207\\
& Introduction to Statistics\\
\end{tabular}
}
\setlength{\baselineskip}{1.05\baselineskip}
\medskip

\begin{center}
\textbf{\large  Confidence Intervals (II)}
\end{center}

%\newcommand{\HH}{\hspace{20pt}\hphantom{a) } }
%
\newcommand{\twobox}[2]{\begin{pspicture}(-0.3,0.2)(2,1.3)
\psframe(0,0)(0.7,0.7)\rput(0.35,0.35){#1}
\psframe(1,0)(1.7,0.7)\rput(1.35,0.35){#2}
\psline(-0.3,0.9)(-0.3,-0.2)(2,-0.2)(2,0.9)(-0.3,0.9)
\end{pspicture}}
%
\newcommand{\threebox}[3]{\begin{pspicture}(0.2,0.2)(2.5,1.3)
\psframe(0,0)(0.7,0.7)\rput(0.35,0.35){#1}
\psframe(1,0)(1.7,0.7)\rput(1.35,0.35){#2}
\psframe(2,0)(2.7,0.7)\rput(2.35,0.35){#3}
\psline(-0.3,0.9)(-0.3,-0.2)(3,-0.2)(3,0.9)(-0.3,0.9)
\end{pspicture}}

\newcommand{\VS}{\vspace{.3in}}


1. (Hypothetical)  A survey organization seeks to estimate the outcome of an upcoming 
election.  They conduct several independent simple random samples of registered voters.
For each sample find a 95\% confidence interval on the percentage of registered
voters who intend to vote for the Coffee Party candidate.

\hspace{20pt} Sample 1:  15 Coffee Party voters out of 100 voters polled.
\vspace{2.2in}

\hspace{20pt} Sample 2:  30 Coffee Party voters out of 225 voters polled.
\vspace{2.2in}

\hspace{20pt} Sample 3:  100 Coffee Party voters out of 1000 voters polled.

\vfill
\eject


2. A box contains tickets with 1s and tickets with 0s.  
Suppose that the fraction of tickets with 1s is 50\%.
What is the chance that a random sample of 100 tickets having 40\% 1s?
\vspace{1.5in}



3. Repeat \#2 if the sample size is 1000.
\vspace{1.5in}

4. Repeat \#2 if the sample size is 10,000.
\vspace{1.5in}

5. Four hundred draws are made at random with replacement from a box
of numbered tickets;  150 are positive.  
Find a 95\% confidence interval
for the percent of positive tickets in the box.
Someone tells you that
50\% of the tickets in the box show positive numbers.  Do you believe it?  
Justify your answer.


\vfill
\eject

\end{document}


